\documentclass[11pt]{article}
% Shared preamble for the hanzo-kernel Paper 07 series.
% One style, one vocabulary, inputted by every paper so the series reads as one voice.
\usepackage[margin=1in]{geometry}
\usepackage{booktabs}
\usepackage{amsmath}
\usepackage{amssymb}
\usepackage{graphicx}
\usepackage{xcolor}
\usepackage{hyperref}
\hypersetup{colorlinks=true, linkcolor=blue, urlcolor=blue, citecolor=blue}
\usepackage{enumitem}
\usepackage{listings}
\usepackage{array}

\definecolor{kw}{rgb}{0.13,0.29,0.53}
\definecolor{cmt}{rgb}{0.40,0.40,0.40}
\definecolor{str}{rgb}{0.16,0.50,0.16}
\lstset{
  basicstyle=\ttfamily\footnotesize,
  keywordstyle=\color{kw}\bfseries,
  commentstyle=\color{cmt}\itshape,
  stringstyle=\color{str},
  showstringspaces=false,
  breaklines=true,
  columns=fullflexible,
  frame=single,
  framesep=4pt,
}
\lstdefinelanguage{Rust}{
  morekeywords={fn,let,mut,pub,struct,impl,trait,enum,match,if,else,for,while,loop,return,use,mod,crate,self,super,async,await,where,type,const,static,unsafe,ref,move,dyn,as,true,false},
  morecomment=[l]{//}, morecomment=[s]{/*}{*/}, morestring=[b]", sensitive=true,
}

\newcommand{\x}{\ensuremath{\times}}
\newcommand{\dpa}{\texttt{dp4a}}
\newcommand{\hk}{\texttt{hanzo-kernel}}
\newcommand{\code}[1]{\texttt{#1}}

% Absorb minor prose overfulls without rivers; keep line-breaking sane.
\tolerance=800
\emergencystretch=3em


\title{\textbf{Intrinsic Islands:\\ An Oracle-Anchored Escape Hatch for a Portable Kernel}}
\author{Hanzo AI --- ML Systems}
\date{hanzo-kernel Paper 07.3 --- part of the \emph{One Source, Every Backend} series}

\begin{document}
\maketitle

\begin{abstract}
Every portable kernel DSL eventually meets a loop that wants a target-specific instruction: the hardware
4-wide int8 dot on one backend, a subgroup shuffle on another. The tempting response is to fork a whole
per-backend kernel for that one line --- which rebuilds, in miniature, the very Cartesian product the DSL
exists to collapse. This paper presents a disciplined alternative: the \emph{intrinsic island}, a scoped,
target-gated block written \emph{inside} the one kernel. The kernel keeps one signature, one launch
contract, one reference, and one test; only the marked inner expression differs per target. The design
rests on two rules. First, selection is a \emph{compile-time} \code{match} on a target tag the launch
layer fills in and the caller never sees, so exactly one arm is lowered per backend and nothing else
compiles. Second, the \code{default} arm is \emph{normative}: it \emph{is} the reference semantics, the
CPU always takes it, and every accelerated arm must equal it byte-for-byte on the CPU oracle before it can
ship. An island is thus an escape hatch that cannot escape correctness --- the accelerated instruction is
never trusted, only proven equivalent. We show the mechanism, the one-kernel-one-oracle test, and why an
island beats a fork on exactly the axis that matters.
\end{abstract}

\section{The pain: one line wants a special instruction}
\label{sec:pain}

A quantized matvec spends its inner loop on a 4-wide integer dot product: four \code{int8} weights times
four \code{int8} activations, summed to an \code{i32}. On NVIDIA, AMD, Apple, and Vulkan hardware this is
a \emph{single instruction} (\dpa{} / \code{OpSDotAccSat}) --- an order of magnitude cheaper than four
scalar multiply-adds. A portable kernel that emits four scalar MACs leaves that instruction on the table;
one that emits the hardware dot cannot run on a target without it (the CPU reference has no such op).

The reflex is to fork: a hand-written \code{matvec\_cuda} that uses the intrinsic, and a portable
\code{matvec\_fallback} that does not. But a fork for one inner line drags the entire kernel with it ---
its signature, its launch bounds, its indexing, its reference, its test --- all now duplicated and free to
drift. That is the Cartesian product (backends $\times$ ops) reappearing at the granularity of a single
expression. We want the hardware instruction where it exists and the portable arithmetic where it does
not, \emph{without} a second copy of the kernel.

\section{The idea: gate the expression, not the kernel}
\label{sec:idea}

An island moves the fork inward, from the whole kernel to the one expression that actually differs
(Listing~\ref{lst:island}). The kernel is written once. At the one spot that needs a target-specific
idiom, an \code{island!} block lists accelerated arms by backend and a mandatory \code{default} arm. The
rest of the kernel --- the loop, the indexing, the accumulation, the bounds --- is shared verbatim.

\begin{lstlisting}[language=Rust,caption={An island in a \dpa{} matvec. The accelerated arm issues the
hardware 4-wide int8 dot; the \code{default} arm --- the normative oracle, taken by the CPU and any
unlisted target --- computes the identical contraction with portable scalar lane MACs. Same math, two
idioms.},label={lst:island}]
let dp = island! {
    // Accelerated: the hardware 4-wide int8 dot (dp4a / OpSDot), one instruction.
    cuda | rocm | metal | vulkan => { w.dot(x) }
    // NORMATIVE oracle: the identical integer contraction, portable scalar lane MACs.
    default => { w[0]*x[0] + w[1]*x[1] + w[2]*x[2] + w[3]*x[3] }
};
\end{lstlisting}

Two properties make this safe rather than merely convenient, and both are structural.

\paragraph{The branch site names backends, not a target.} The author writes \code{cuda | rocm | metal |
vulkan => ...}; they never write ``if I am on CUDA.'' The \emph{tag} that selects the arm is filled in
elsewhere (\S\ref{sec:mech}), so a kernel author cannot accidentally couple the math to a runtime string.

\paragraph{\code{default} is the definition.} The accelerated arms are optimizations \emph{of} the
\code{default} arm. \code{default} is not a fallback in the sense of ``lower quality if unlucky'' --- it
is the \emph{normative} semantics that the accelerated arms must reproduce. This is the inversion that
makes an island trustworthy: the portable arithmetic is the truth, and the fast instruction has to match
it.

\section{The mechanism: target selection at compile time}
\label{sec:mech}

The subtlety is that the lowering engine compiles the same target-agnostic IR for every backend; at
kernel-expansion time there is no ``which backend am I'' string available inside the kernel, only a
capability set. But the backend brand \emph{is} known on the \emph{host}, from the runtime's name
(\code{"cpu"}, \code{"cuda"}, \code{"hip"}, \code{"wgpu<spirv>"}, \code{"wgpu<msl>"}). So an island is
selected by a compile-time tag threaded in from the host (Listing~\ref{lst:target}):

\begin{lstlisting}[language=Rust,caption={The target tag is derived on the host from the runtime name and
threaded into the kernel as a \code{\#[comptime]} parameter. The \code{\#[kernel]} macro rewrites
\code{island!\{...\}} into a comptime \code{match target}. Unknown names fall to \code{Cpu} --- fail
safe, never a wrong accelerated arm.},label={lst:target}]
pub enum Target { Cpu, Cuda, Rocm, Metal, Vulkan, WebGpu }   // #[comptime] tag: Eq + Hash + Copy

pub fn from_runtime_name(name: &str) -> Target {
    match name {
        "cuda" => Target::Cuda,
        "hip"  => Target::Rocm,
        n if n.contains("spirv") => Target::Vulkan,
        n if n.contains("msl")   => Target::Metal,
        n if n.starts_with("wgpu") => Target::WebGpu,
        _ => Target::Cpu,                                    // fail safe -> the oracle arm
    }
}
\end{lstlisting}

Because \code{target} is a \code{\#[comptime]} value, the \code{match} it drives is a comptime match over
a constant scrutinee: \textbf{only the selected arm is lowered.} Every backend compiles exactly its own
arm and nothing else --- there is no runtime branch, no dead code for other targets in the shipped kernel.
The tag is kept out of the authoring surface twice over: the branch site lists backends (not a target),
and the host launch wrapper fills the tag (\code{Target::of::<R>}), so callers pass only data. One
signature, one launch, every backend.

\section{The oracle contract, tested}
\label{sec:oracle}

Here is the rule that makes an island an escape hatch that cannot escape correctness. The CPU runtime is
\emph{always} tagged \code{Target::Cpu} --- a variant no accelerated arm ever claims --- so on the CPU
every island resolves to \code{default}. The CPU runtime is therefore the single bit-exact oracle for
every island (the general form of this law, and its cross-backend enforcement, is the companion oracle
paper, \cite{oracle}). The test states the contract directly and mechanically: run the kernel on the CPU
with each accelerated tag \emph{forced}, and it must equal the \code{default} result byte-for-byte
(Listing~\ref{lst:oracletest}).

\begin{lstlisting}[language=Rust,caption={The oracle-equivalence test. Every accelerated arm, executed on
the CPU runtime with its tag forced, must be byte-identical to the \code{default} arm and to the
plain-Rust reference. An accelerated instruction is never trusted --- it is proven equal to the normative
arm.},label={lst:oracletest}]
// The default arm (CPU is tagged Cpu -> default) must equal the plain-Rust oracle.
let default = matvec_island_run_with::<CpuRuntime>(&c, &wq, &xq, &wd, rows, k, Target::Cpu);
assert_eq!(bits(&default), bits(&matvec_island_ref(&wq, &xq, &wd, rows, k)));

// Every accelerated arm, forced onto the CPU runtime, must equal the oracle byte-for-byte.
for accel in [Target::Cuda, Target::Rocm, Target::Metal, Target::Vulkan, Target::WebGpu] {
    let out = matvec_island_run_with::<CpuRuntime>(&c, &wq, &xq, &wd, rows, k, accel);
    assert_eq!(bits(&out), bits(&default), "island arm {accel:?} diverged from default");
}
\end{lstlisting}

A separate test asserts the \emph{production} launch --- the one that derives the tag from the runtime and
hides it --- resolves to the oracle arm on the CPU without the caller naming a target. So the property is
proven both under forced tags (every arm equals \code{default}) and under the real launch path (the
caller-facing API does the right thing by default).

\section{Why an island beats a fork}
\label{sec:vsfork}

The comparison is not aesthetic; it is about what can drift and what can be verified
(Table~\ref{tab:vsfork}). A fork gives two kernels: two signatures to keep in sync, two references to
trust, two tests to keep honest, and no mechanical guarantee that the fast kernel computes what the
portable one does --- the equivalence lives only in a human's intention. An island gives one kernel: one
signature, one reference (the \code{default} arm), one test that \emph{proves} every accelerated arm
equals it. The accelerated instruction is contained, not trusted.

\begin{table}[h]
\centering\small
\setlength{\tabcolsep}{7pt}
\begin{tabular}{@{}lll@{}}
\toprule
 & \textbf{Fork (per-backend kernel)} & \textbf{Island (scoped block)} \\
\midrule
Signatures     & one per backend, drift-prone      & one, shared \\
Reference      & one per fork, trust-based         & one \code{default} arm, normative \\
Equivalence    & human intention only              & proven byte-for-byte on the CPU oracle \\
What differs   & the whole kernel                  & the one gated expression \\
New backend    & another full kernel               & another arm, or falls to \code{default} \\
\bottomrule
\end{tabular}
\caption{An island localizes the difference to the expression that actually differs and anchors it to a
CPU-checkable normative arm. A fork localizes nothing and anchors nothing.}
\label{tab:vsfork}
\end{table}

\section{The designed seam}

An island is not a licence to sprinkle intrinsics; it is the \emph{one} sanctioned place a per-backend
idiom may live, precisely because it comes with the oracle contract attached. That matters most for the
operation this series shows does \emph{not} collapse into portable form: block-quantized tensor-core MMQ,
whose viable implementation is per-backend by nature (\cite{mmq}). Even there, an island holds the
hand-written tensor-core arm behind a normative, CPU-checkable \code{default}, so the one un-collapsed
kernel still keeps one signature, one oracle, one test. The escape hatch exists so that the rare
irreducibly-per-backend kernel does not force a return to forks --- it is the boundary condition of the
whole ``one source, every backend'' program, made safe.

\begin{thebibliography}{9}
\small
\bibitem{umbrella} Hanzo AI. \emph{One Source, Every Backend: The hanzo-kernel DSL.} Paper 07 (umbrella).
\bibitem{fuse} Hanzo AI. \emph{Fusion Is Composition: Kernel Fusion as the Functor Law.} Paper 07.1.
\bibitem{tune} Hanzo AI. \emph{The Schedule Is a Value: Comptime-Specialized Autotuning.} Paper 07.2.
\bibitem{island} Hanzo AI. \emph{Intrinsic Islands: An Oracle-Anchored Escape Hatch.} Paper 07.3 (this paper).
\bibitem{oracle} Hanzo AI. \emph{One Oracle, Every Backend: The CPU Bit-Exactness Law.} Paper 07.4.
\bibitem{mmq} Hanzo AI. \emph{When the Product Won't Collapse: An int8 Tensor-Core MMQ Negative Result.} Paper 07.5.
\bibitem{method} Hanzo AI. \emph{Verify, Then Delete: Collapsing a Kernel Cartesian Product.} Paper 07.6.
\end{thebibliography}

\end{document}
